\documentclass{article}
\usepackage{import}
\import{../../template/}{article-template}
\graphicspath{{../../images/}}
\addbibresource{../../template/library.bib}

\usepackage{tipa}

\title{English}
\author{PENG}
\date{Created: 20260913, Version: 20260913}

\begin{document}
\maketitle
\tableofcontents
\section{Vocabulary}
\url{https://dictionary.cambridge.org/}

\subsection{adjective}
\begin{itemize}
    \item noun \textipa{["\ae dZ.ek.tIv]}
          \begin{itemize}
              \item a word that describes a noun or pronoun
          \end{itemize}
\end{itemize}

\subsection{euphoric}
\begin{itemize}
    \item adjective \textipa{[ju:"f\textturnscripta r.Ik]}
          \begin{itemize}
              \item extremely happy and excited
                    \begin{itemize}
                        \item a euphoric mood
                    \end{itemize}
          \end{itemize}
\end{itemize}

\subsection{epiphany}
\begin{itemize}
    \item noun \textipa{[I"pIf.\textsuperscript{@}n.i]}
          \begin{itemize}
              \item a moment when you suddenly feel that you understand, or suddenly become conscious of, something that is very important to you
                    \begin{itemize}
                        \item Euphoric, mind-bending \href{https://www.engineering.upenn.edu/~cis1940/spring13/}{epiphanies} may result!
                    \end{itemize}
          \end{itemize}
\end{itemize}

\subsection{execute}
\begin{itemize}
    \item verb \textipa{["ek.sI.kju:t]}
          \begin{itemize}
              \item to do or perform something, especially in a planned way
                    \begin{itemize}
                        \item \href{https://www.engineering.upenn.edu/~cis1940/spring13/lectures/01-intro.html}{execute} instructions
                    \end{itemize}
          \end{itemize}
\end{itemize}

\subsection{mutation}
\begin{itemize}
    \item noun \textipa{[mju:"teI.S@n]}
          \begin{itemize}
              \item the way in which genes change and produce permanent differences
                    \begin{itemize}
                        \item No \href{https://www.engineering.upenn.edu/~cis1940/spring13/lectures/01-intro.html}{mutation}! Everything (variables, data structure) is immutable.
                    \end{itemize}
          \end{itemize}
\end{itemize}

\subsection{imperative}
\begin{itemize}
    \item adjective \textipa{[Im"per.@.tIv]}
          \begin{itemize}
              \item extremely important or urgent
              \item used for giving an instruction or order
                    \begin{itemize}
                        \item an \href{https://www.engineering.upenn.edu/~cis1940/spring13/lectures/01-intro.html}{imperative} or object-oriented paradigm
                    \end{itemize}
          \end{itemize}
\end{itemize}

\subsection{immense}
\begin{itemize}
    \item adjective \textipa{[i"mens]}
          \begin{itemize}
              \item extremely large in size or degree
                    \begin{itemize}
                        \item \href{https://www.engineering.upenn.edu/~cis1940/spring13/lectures/01-intro.html}{immense} help
                    \end{itemize}
          \end{itemize}
\end{itemize}

\subsection{facetious}
\begin{itemize}
    \item adjective \textipa{[f@"si:.S@s]}
          \begin{itemize}
              \item not serious about a serious subject, in an attempt to be funny or to appear clever
                    \begin{itemize}
                        \item ``If it compiles, it must be correct'' is mostly \href{https://www.engineering.upenn.edu/~cis1940/spring13/lectures/01-intro.html}{facetious}.
                    \end{itemize}
          \end{itemize}
\end{itemize}

\subsection{mantra}
\begin{itemize}
    \item noun \textipa{["m\ae n.tr@]}
          \begin{itemize}
              \item a word or phrase that is often repeated and expresses a particular strong believe
                    \begin{itemize}
                        \item ``Don't Repeat Yourself'' is a \href{https://www.engineering.upenn.edu/~cis1940/spring13/lectures/01-intro.html}{mantra} often heard in the world of programming.
                    \end{itemize}
          \end{itemize}
\end{itemize}

\subsection{auxiliary}
\begin{itemize}
    \item adjective \textipa{[O:g"zIl.i.@.ri]}
          \begin{itemize}
              \item giving help or support, especially to a more important person or thing
                    \begin{itemize}
                        \item auxiliary boundary conditions \cite{EVANS_2022_PARTIAL_DIFFERENTIAL_EQUATIONS}
                    \end{itemize}
          \end{itemize}
\end{itemize}

\subsection{permutation}
\begin{itemize}
    \item noun \textipa{[""p\textrevepsilon :.mju:"teI.S@n]}
          \begin{itemize}
              \item any of the various ways in which a set of things can be ordered
                    \begin{itemize}
                        \item This sorting function returns a \href{https://lean-lang.org/functional_programming_in_lean/Getting-to-Know-Lean/Types/#getting-to-know-types}{permutation} of its input.
                    \end{itemize}
          \end{itemize}
\end{itemize}

\subsection{bureaucratic}
\begin{itemize}
    \item adjective \textipa{[""bjU@.r@"kr\ae t.Ik]}
          \begin{itemize}
              \item involving complicated rules and processes that make something slow and difficult
                    \begin{itemize}
                        \item The program become long and bureaucratic.
                    \end{itemize}
          \end{itemize}
\end{itemize}

\subsection{conundrum}
\begin{itemize}
    \item noun \textipa{[k@"n\textturnv n.dr@m]}
          \begin{itemize}
              \item a question that is a trick, often involving a humorous use of words that have two meanings
                    \begin{itemize}
                        \item \verb|Nat.zero| is a bit of a \href{https://lean-lang.org/functional_programming_in_lean/Getting-to-Know-Lean/Datatypes-and-Patterns/#datatypes-and-patterns}{conumdrum}.
                    \end{itemize}
          \end{itemize}
\end{itemize}

\subsection{polymorphism}
\begin{itemize}
    \item noun \textipa{[""p\textturnscripta l.i"mO:.fI.z@m]}
          \begin{itemize}
              \item the fact that something such as an animal or organism can exist in different forms
                    \begin{itemize}
                        \item The term \href{https://lean-lang.org/functional_programming_in_lean/Getting-to-Know-Lean/Polymorphism/#polymorphism}{polymorphism} typically refers to datatypes and definitions that take types as arguments.
                    \end{itemize}
          \end{itemize}
\end{itemize}

\subsection{customary}
\begin{itemize}
    \item adjective \textipa{["k\textturnv s.t@.m@r.i]}
          \begin{itemize}
              \item traditional
                    \begin{itemize}
                        \item It is \href{https://lean-lang.org/functional_programming_in_lean/Getting-to-Know-Lean/Polymorphism/#polymorphism}{customary} to use Greek letters to name type arguments in Lean.
                    \end{itemize}
          \end{itemize}
\end{itemize}

\subsection{reconcile}
\begin{itemize}
    \item verb \textipa{["rek.@n.saIl]}
          \begin{itemize}
              \item to find a way in which two situations or beliefs that are opposed to each other can agree and exist together
                    \begin{itemize}
                        \item Programs are difficult to \href{https://lean-lang.org/functional_programming_in_lean/Getting-to-Know-Lean/Polymorphism/#polymorphism}{reconcile} with the understanding of computation as the evaluation of mathematical expressions.
                    \end{itemize}
          \end{itemize}
\end{itemize}

\subsection{tedious}
\begin{itemize}
    \item adjective \textipa{["ti:.di.@s]}
          \begin{itemize}
              \item boring
                    \begin{itemize}
                        \item Most proofs are written for people to understand, and leave out many \href{https://lean-lang.org/functional_programming_in_lean/Interlude___-Propositions___-Proofs___-and-Indexing/#props-proofs-indexing}{tedious} details.
                    \end{itemize}
          \end{itemize}
\end{itemize}

\subsection{succinct}
\begin{itemize}
    \item adjective \textipa{[s@k"sINkt]}
          \begin{itemize}
              \item said in a clear and short way; expressing what needs to be said without unnecessary words
                    \begin{itemize}
                        \item A reference work in which all aspects of Lean are clearly specified, and demonstrated through \href{https://lean-lang.org/learn/}{succinct} examples.
                    \end{itemize}
          \end{itemize}
\end{itemize}

\printbibliography
\end{document}
